\section{Экспериментальные исследования и оценка эффективности}
\subsection{Методика тестирования}
Экспериментальная оценка проводилась на тестовом стенде, имитирующем промышленные условия эксплуатации. Аппаратная конфигурация включала сервер с двумя процессорами Intel Xeon Gold 6348 (28 ядер, 56 потоков каждый), 512 ГБ DDR4-3200, NVMe-накопитель для логирования и сетевую карту Mellanox ConnectX-5 с поддержкой 25 Гбит/с. Программное обеспечение: Ubuntu 22.04 LTS, ядро 5.15, компилятор GCC 11.3 с флагами оптимизации \texttt{-O3 -march=native -flto}. Для генерации тестового трафика использовались датасеты CIC-IDS2017 и MAWI Archive, содержащие смешанный легитимный и аномальный трафик с известными векторами атак \cite{sharafaldin2018cic}.

Методика включала три этапа: (1) измерение пропускной способности и задержки при линейном увеличении числа генераторов событий; (2) оценка масштабируемости при изменении числа worker-потоков и конфигурации очередей; (3) тестирование механизма горячей перезагрузки обработчиков и устойчивости к пиковым нагрузкам. В качестве базовых систем для сравнения использовались синтетические реализации на базе Boost.Signal2 и Zeek-подобного скриптового движка, настроенные на идентичную аппаратную платформу и эквивалентную логику обработки.

\subsection{Метрики и результаты}
Пропускная способность измерялась в миллионах событий в секунду (Meps), задержка — в микросекундах (мкс) с распределением по перцентилям (p50, p95, p99). При нагрузке 10 Гбит/с и конфигурации 16 worker-потоков разработанная система продемонстрировала среднюю пропускную способность 2.14 Meps, что на 34\% превышает показатели синтетического аналога на Boost.Signal2 и на 28\% — Zeek-подобной реализации. Задержка диспетчеризации составила p50=38 мкс, p95=82 мкс, p99=147 мкс, что находится в пределах допустимых значений для систем реального времени.

Потребление памяти демонстрировало линейную зависимость от числа активных типов событий и глубины очередей. При максимальной нагрузке пиковое потребление составило 412 МБ, при этом фрагментация памяти не превышала 2.3\%, что подтверждает эффективность пулового аллокатора. Тестирование механизма горячей перезагрузки обработчиков показало среднее время переключения конфигурации 140 мс без потери пакетов и сброса состояния очередей, что соответствует требованиям SLA для высокодоступных инфраструктур \cite{sre2018google}.

\subsection{Анализ масштабируемости и устойчивости}
Исследование масштабируемости проводилось в диапазоне от 1 до 56 потоков. Результаты подтвердили линейный рост пропускной способности вплоть до 32 потоков, после чего наблюдалось плато, обусловленное ограничением пропускной способности межпроцессорной шины (QPI/UPI) и contention на уровне сетевых очередей. Применение work-stealing и NUMA-aware распределения памяти позволило сократить деградацию производительности на узких местах до 8\%. Устойчивость к сбоям оценивалась путём внедрения искусственных задержек и исключений в обработчики. Благодаря изоляции контекстов выполнения и механизму backpressure система сохранила работоспособность, сбросив 0.4\% низкоприоритетных событий при пиковой перегрузке, что предотвратило каскадный отказ конвейера \cite{dag2022design}.

Сравнительный анализ показал, что разработанный компонент обеспечивает предсказуемое поведение в условиях высокой нагрузки, снижает время интеграции новых аналитических модулей за счёт декларативного API и устраняет необходимость модификации ядра DPI-систем при обновлении логики обработки. Полученные метрики подтверждают применимость системы в распределённых NTA-инфраструктурах, требующих высокой точности детектирования и минимального времени реакции на изменения сетевой среды.

\newpage